% !TEX root = ../main.tex
\chapter{Project Management - Liam}

\section{Planung und Organisation}
Die organisatorische Planung und Umsetzung des Projekts \glqq Ascenta\grqq\ basierte auf einer meilensteinorientierten, agilen Vorgehensweise. Zu Beginn des Projekts (Vorbereitungsphase) lag der Fokus auf der grundlegenden organisatorischen Planung, der Definition von Grob- und Feinzielen sowie der Recherche und Beschaffung der elektronischen Komponenten. 

Die formelle Strukturierung der Arbeitsschritte und die Einreichung der Projektmeilensteine erfolgten über das schulspezifische ABA-Portal. Um eine strukturierte Arbeitsteilung zwischen den Bereichen Software, Elektronik und 3D-Druck zu gewährleisten, wurde der Projektplan in wöchentliche Sprints (Kalenderwochen, KW) unterteilt. Zur Überwachung des Fortschritts und der individuellen Arbeitslast wurde eine detaillierte Zeiterfassung (Timesheet) mittels Tabellenkalkulation implementiert, welche die geplanten Freizeitstunden dem tatsächlichen Aufwand (Soll/Ist-Vergleich) kontinuierlich gegenüberstellt.

\section{Evaluierung der Arbeitsstunden}
Die wöchentliche Evaluierung der Tätigkeitsberichte ermöglichte eine präzise Überwachung des Projektfortschritts. Ein Abgleich der geplanten Ressourcen mit dem tatsächlichen Aufwand zeigt insbesondere in der initialen Projektphase eine deutliche Diskrepanz: Während für die Projektdefinition anfänglich 6 Stunden pro Teammitglied veranschlagt wurden, belief sich der tatsächliche Aufwand auf 20 Stunden pro Person. Diese anfängliche Unterschätzung ist typisch für komplexe Hardware-Software-Kombinationen, da grundlegende Architektur-Entscheidungen zeitintensiv sind.

In den darauffolgenden Wochen (beispielsweise in KW 37 und 38) stabilisierte sich die Arbeitslast und entsprach mit durchschnittlich 6 Stunden pro Teammitglied exakt den Planwerten. Die konsequente Dokumentation der wöchentlichen Ist-Zustände – wie etwa der erfolgreiche Test des Spannungsbetriebsbereichs der Akkus oder die Evaluation der Micro-Speech-Software – half maßgeblich dabei, technische Hürden frühzeitig zu identifizieren und den Projektplan dynamisch anzupassen. 

\section{Zeitplan}
Die nachfolgenden Tabellen dokumentieren die individuell erbrachten Arbeitsstunden für die Diplomarbeit, unterteilt in reguläre Schulstunden (7,5\,h pro Woche; bereinigt um Ferien und Krankenstände) und zusätzliche Freizeitstunden.

%LIAM
\renewcommand{\arraystretch}{1.1}
\begin{longtable}{l r r p{\dimexpr\textwidth-4.5cm\relax}}
\caption{Wöchentliche Zeiterfassung - Liam}
\label{tab:zeiterfassung_liam} \\
\hline
\textbf{KW} & \multicolumn{2}{c}{\textbf{Arbeitszeit [h]}} & \textbf{Aufgabenbeschreibung} \\
& Schule & Freizeit & \\
\hline
\endfirsthead
\caption[]{Wöchentliche Zeiterfassung - Liam (Fortsetzung)} \\
\hline
\textbf{KW} & \multicolumn{2}{c}{\textbf{Arbeitszeit [h]}} & \textbf{Aufgabenbeschreibung} \\
& Schule & Freizeit & \\
\hline
\endhead
\hline
\multicolumn{4}{r}{\textit{Fortsetzung auf der nächsten Seite...}} \\
\endfoot
\hline
\textbf{Summe} & \textbf{165,0} & \textbf{237,0} & \textbf{Gesamt: 402,0 h} \\
\hline
\hline
\endlastfoot
Vorb. & 0,0 & 20,0 & Elektronische Komponenten bestellt. Organisatorische Planung \& Grobe Ziele. \\
37 & 7,5 & 6,0 & ABA Portal begutachtet, Audio TTS via Eingabe und Konsolen-Ausgabe. \\
38 & 7,5 & 6,0 & ABA Portal eingereicht. Akkus getestet. Micro-Speech Software evaluiert. \\
39 & 7,5 & 10,0 & Bügel-Prototyp gedruckt. 3D-Drucker konfiguriert. Spannungsteiler aufgebaut. \\
40 & 7,5 & 6,0 & KiCad Platinendesign (Spannungsteiler). Brillenhalterung. Modell trainiert. \\
41 & 7,5 & 10,0 & BLE-Verbindung App $\leftrightarrow$ Nicla. Brillendesign/Bügel \& KiCad angepasst. \\
42 & 7,5 & 5,0 & Brillendesign-Fehler behoben. Bilder erfolgreich übertragen. Button Interrupt. \\
43 & 7,5 & 12,0 & Scharniere/Halterungen modelliert. Interrupt Code mit Nicla zusammengeführt. \\
44 & 0,0 & 3,0 & Einrastfunktion modelliert. STT über App. Code Verbesserung Interrupt. \\
45 & 7,5 & 6,0 & Plakat, PR. Fuel-Gauge Werte an Webapp dargestellt. 3D-Druck Prototyp\_2. \\
46 & 7,5 & 12,0 & Audio aufgenommen \& an App gesendet. Wettbewerb JugendInnovativ geplant. \\
47 & 7,5 & 10,0 & Audio-Erfolg mit wit.ai STT \& Handy-TTS. Start EdgeImpulse Objekterkennung. \\
48 & 7,5 & 8,0 & App-Prototyp (Foto/Audio-UI). Objekterkennung Versuche. Nasenbügel. \\
49 & 7,5 & 15,0 & Objekterkennung Bugs fixen. Letzte 3D-Druck Anpassungen für den TdoT. \\
50 & 7,5 & 6,0 & Objekterkennung neue Gegenstände. 3D-Druck: Fehler ausgemacht \& skizziert. \\
51 & 7,5 & 6,0 & Objekterkennung AI Boundingboxes (owvl2). Neue Bügel für 620\,mAh Akkus. \\
52 & 0,0 & 2,0 & Evaluierung der Objekterkennung (FOMO-Architektur). \\
01 & 0,0 & 3,0 & Objekterkennung FOMO. Bügel-Design für 620\,mAh Akkus in 3D modelliert. \\
02 & 0,0 & 2,0 & Halterung mit neuen Scharnieren. Planung OE in App. YOLO Testung am PC. \\
03 & 0,0 & 0,0 & Technische Schwierigkeiten \& Debugging. \\
04 & 7,5 & 6,0 & Fast vollständiger Projektbericht für Jugend Innovativ erstellt. \\
05 & 7,5 & 15,0 & JI-Bericht abgegeben. Lokale STT erfolgreich am Handy implementiert. \\
06 & 7,5 & 6,0 & Diplomarbeit Struktur ausarbeiten und mit Schreiben beginnen. \\
07 & 7,5 & 15,0 & Planung Umstieg auf WiFi. IMU als Input-Trigger. Edge Impulse Model. \\
08 & 0,0 & 20,0 & Erfolgreiche Implementierung der TCP/WiFi-Datenübertragung. \\
09 & 7,5 & 10,0 & Finale Software auf GitHub. Demo (Objekterkennung/Audio-Tap). Bügel optimiert. \\
10 & 7,5 & 6,0 & 3D-Druck: Schalter-Implementierung. Diplomarbeit verfasst. \\
11 & 7,5 & 8,0 & Content geschrieben (Related Work, Abstract, Future Work, Konzept, Umsetzung). \\
12 & 7,5 & 3,0 & Diplomarbeit layouten und fertigstellen. \\
\end{longtable}
\newpage

%FABIAN

\renewcommand{\arraystretch}{1.1}
\begin{longtable}{l r r p{\dimexpr\textwidth-4.5cm\relax}}
\caption{Wöchentliche Zeiterfassung - Fabian}
\label{tab:zeiterfassung_fabian} \\
\hline
\textbf{KW} & \multicolumn{2}{c}{\textbf{Arbeitszeit [h]}} & \textbf{Aufgabenbeschreibung} \\
& Schule & Freizeit & \\
\hline
\endfirsthead

\caption[]{Wöchentliche Zeiterfassung - Fabian (Fortsetzung)} \\
\hline
\textbf{KW} & \multicolumn{2}{c}{\textbf{Arbeitszeit [h]}} & \textbf{Aufgabenbeschreibung} \\
& Schule & Freizeit & \\
\hline
\endhead

\hline
\multicolumn{4}{r}{\textit{Fortsetzung auf der nächsten Seite...}} \\
\endfoot

\hline
\textbf{Summe} & \textbf{165,0} & \textbf{180,0} & \textbf{Gesamt: 345,0 h} \\
\hline
\hline
\endlastfoot
Vorb. & 0,0 & 20,0 & Elektronische Komponenten bestellt. Organisatorische Planung \& Grobe Ziele. \\
37 & 7,5 & 6,0 & ABA Portal begutachtet, Audio TTS via Eingabe und Konsolen-Ausgabe. \\
38 & 7,5 & 6,0 & ABA Portal eingereicht. Akkus getestet. Micro-Speech Software evaluiert. \\
39 & 7,5 & 6,0 & Bügel-Prototyp gedruckt. 3D-Drucker konfiguriert. Spannungsteiler aufgebaut. \\
40 & 7,5 & 6,0 & KiCad Platinendesign (Spannungsteiler). Brillenhalterung. Modell trainiert. \\
41 & 7,5 & 6,0 & BLE-Verbindung App $\leftrightarrow$ Nicla. Brillendesign/Bügel \& KiCad angepasst. \\
42 & 7,5 & 3,0 & Brillendesign-Fehler behoben. Bilder erfolgreich übertragen. Button Interrupt. \\
43 & 7,5 & 12,0 & Scharniere/Halterungen modelliert. Interrupt Code mit Nicla zusammengeführt. \\
44 & 0,0 & 10,0 & Einrastfunktion modelliert. STT über App. Code Verbesserung Interrupt. \\
45 & 7,5 & 6,0 & Plakat, PR. Fuel-Gauge Werte an Webapp dargestellt. 3D-Druck Prototyp\_2. \\
46 & 7,5 & 4,0 & Audio aufgenommen \& an App gesendet. Wettbewerb JugendInnovativ geplant. \\
47 & 7,5 & 6,0 & Audio-Erfolg mit wit.ai STT \& Handy-TTS. Start EdgeImpulse Objekterkennung. \\
48 & 7,5 & 6,0 & App-Prototyp (Foto/Audio-UI). Objekterkennung Versuche. Nasenbügel. \\
49 & 7,5 & 6,0 & Objekterkennung Bugs fixen. Letzte 3D-Druck Anpassungen für den TdoT. \\
50 & 7,5 & 3,0 & Objekterkennung neue Gegenstände. 3D-Druck: Fehler ausgemacht \& skizziert. \\
51 & 7,5 & 10,0 & Objekterkennung AI Boundingboxes (owvl2). Neue Bügel für 620\,mAh Akkus. \\
52 & 0,0 & 6,0 & Evaluierung der Objekterkennung (FOMO-Architektur). \\
01 & 0,0 & 6,0 & Objekterkennung FOMO. Bügel-Design für 620\,mAh Akkus in 3D modelliert. \\
02 & 7,5 & 4,0 & Halterung mit neuen Scharnieren. Planung OE in App. YOLO Testung am PC. \\
03 & 0,0 & 0,0 & Technische Schwierigkeiten \& Debugging. \\
04 & 0,0 & 0,0 & Vorarbeiten für Jugend Innovativ Projektbericht. \\
05 & 7,5 & 6,0 & JI-Bericht abgegeben. Lokale STT erfolgreich am Handy implementiert. \\
06 & 7,5 & 6,0 & Diplomarbeit Struktur ausarbeiten und mit Schreiben beginnen. \\
07 & 7,5 & 6,0 & Planung Umstieg auf WiFi. IMU als Input-Trigger. Edge Impulse Model. \\
08 & 0,0 & 6,0 & Erfolgreiche Implementierung der TCP/WiFi-Datenübertragung. \\
09 & 7,5 & 6,0 & Finale Software auf GitHub. Demo (Objekterkennung/Audio-Tap). Bügel optimiert. \\
10 & 7,5 & 6,0 & 3D-Druck: Schalter-Implementierung. Diplomarbeit verfasst. \\
11 & 7,5 & 6,0 & Schreiben der Diplomarbeit (Elektronik, Bauteilbeschreibungen, Probleme). \\
12 & 7,5 & 6,0 & Diplomarbeit layouten und fertigstellen. \\
\end{longtable}
\newpage

%MANUEL
\renewcommand{\arraystretch}{1.1}
\begin{longtable}{l r r p{\dimexpr\textwidth-4.5cm\relax}}
\caption{Wöchentliche Zeiterfassung - Manuel}
\label{tab:zeiterfassung_manuel} \\
\hline
\textbf{KW} & \multicolumn{2}{c}{\textbf{Arbeitszeit [h]}} & \textbf{Aufgabenbeschreibung} \\
& Schule & Freizeit & \\
\hline
\endfirsthead

\caption[]{Wöchentliche Zeiterfassung - Manuel (Fortsetzung)} \\
\hline
\textbf{KW} & \multicolumn{2}{c}{\textbf{Arbeitszeit [h]}} & \textbf{Aufgabenbeschreibung} \\
& Schule & Freizeit & \\
\hline
\endhead

\hline
\multicolumn{4}{r}{\textit{Fortsetzung auf der nächsten Seite...}} \\
\endfoot

\hline
\textbf{Summe} & \textbf{165,0} & \textbf{193,0} & \textbf{Gesamt: 358,0 h} \\
\hline
\hline
\endlastfoot
Vorb. & 0,0 & 20,0 & Elektronische Komponenten bestellt. Organisatorische Planung \& Grobe Ziele. \\
37 & 7,5 & 6,0 & ABA Portal begutachtet, Audio TTS via Eingabe und Konsolen-Ausgabe. \\
38 & 7,5 & 6,0 & ABA Portal eingereicht. Akkus getestet. Micro-Speech Software evaluiert. \\
39 & 7,5 & 15,0 & Bügel-Prototyp gedruckt. 3D-Drucker konfiguriert. Spannungsteiler aufgebaut. \\
40 & 7,5 & 6,0 & KiCad Platinendesign (Spannungsteiler). Brillenhalterung. Modell trainiert. \\
41 & 7,5 & 9,0 & BLE-Verbindung App $\leftrightarrow$ Nicla. Brillendesign/Bügel \& KiCad angepasst. \\
42 & 7,5 & 6,0 & Brillendesign-Fehler behoben. Bilder erfolgreich übertragen. Button Interrupt. \\
43 & 7,5 & 10,0 & Scharniere/Halterungen modelliert. Interrupt Code mit Nicla zusammengeführt. \\
44 & 0,0 & 3,0 & Einrastfunktion modelliert. STT über App. Code Verbesserung Interrupt. \\
45 & 7,5 & 3,0 & Plakat, PR. Fuel-Gauge Werte an Webapp dargestellt. 3D-Druck Prototyp\_2. \\
46 & 7,5 & 10,0 & Audio aufgenommen \& an App gesendet. Wettbewerb JugendInnovativ geplant. \\
47 & 7,5 & 7,0 & Audio-Erfolg mit wit.ai STT \& Handy-TTS. Start EdgeImpulse Objekterkennung. \\
48 & 7,5 & 6,0 & App-Prototyp (Foto/Audio-UI). Objekterkennung Versuche. Nasenbügel. \\
49 & 7,5 & 12,0 & Objekterkennung Bugs fixen. Letzte 3D-Druck Anpassungen für den TdoT. \\
50 & 7,5 & 6,0 & Objekterkennung neue Gegenstände. 3D-Druck: Fehler ausgemacht \& skizziert. \\
51 & 7,5 & 6,0 & Objekterkennung AI Boundingboxes (owvl2). Neue Bügel für 620\,mAh Akkus. \\
52 & 0,0 & 0,0 & Evaluierung der Objekterkennung (FOMO-Architektur). \\
01 & 0,0 & 10,0 & Objekterkennung FOMO. Bügel-Design für 620\,mAh Akkus in 3D modelliert. \\
02 & 7,5 & 4,0 & Halterung mit neuen Scharnieren. Planung OE in App. YOLO Testung am PC. \\
03 & 0,0 & 0,0 & Technische Schwierigkeiten \& Debugging. \\
04 & 7,5 & 6,0 & Fast vollständiger Projektbericht für Jugend Innovativ erstellt. \\
05 & 0,0 & 0,0 & JI-Bericht abgegeben. Lokale STT erfolgreich am Handy implementiert. \\
06 & 7,5 & 4,0 & Diplomarbeit Struktur ausarbeiten und mit Schreiben beginnen. \\
07 & 7,5 & 10,0 & Planung Umstieg auf WiFi. IMU als Input-Trigger. Edge Impulse Model. \\
08 & 0,0 & 6,0 & Erfolgreiche Implementierung der TCP/WiFi-Datenübertragung. \\
09 & 7,5 & 6,0 & Finale Software auf GitHub. Demo (Objekterkennung/Audio-Tap). Bügel optimiert. \\
10 & 7,5 & 6,0 & 3D-Druck: Schalter-Implementierung. Diplomarbeit verfasst. \\
11 & 7,5 & 7,0 & Schreiben der Diplomarbeit (Schalterhalterung \& Konstruktionstext). \\
12 & 7,5 & 3,0 & Diplomarbeit layouten und fertigstellen. \\
\end{longtable}